\documentclass[11pt,a4paper]{article}
\usepackage[utf8]{inputenc}
\usepackage[T1]{fontenc}
\usepackage{lmodern}
\usepackage[french]{babel}
\usepackage{amsmath,amssymb,mathtools}
\usepackage{icomma}
\usepackage{xcolor}
\usepackage{colortbl}
\usepackage{tikz}
\usepackage[most]{tcolorbox}
\usepackage{enumitem}
\usepackage{array}
\usepackage{booktabs}
\usepackage{fancyhdr}
\usepackage[hidelinks]{hyperref}
\usetikzlibrary{arrows.meta,positioning,calc,patterns,backgrounds,shapes.geometric,decorations.pathreplacing,decorations.markings}
\usepackage[a4paper,margin=2.15cm,headheight=26pt,headsep=12pt]{geometry}
\usepackage{titlesec}

% ---------------- Palette ----------------
\definecolor{primary}{RGB}{0,151,169}
\definecolor{primarydark}{RGB}{0,88,101}
\definecolor{defcol}{RGB}{0,114,178}
\definecolor{thmcol}{RGB}{0,140,120}
\definecolor{formulacol}{RGB}{198,93,9}
\definecolor{excol}{RGB}{0,135,90}
\definecolor{attncol}{RGB}{200,40,55}
\definecolor{methcol}{RGB}{90,90,100}
\definecolor{remarkcol}{RGB}{120,94,200}
\definecolor{barcol}{RGB}{0,151,169}
\definecolor{barcold}{RGB}{0,110,124}
\definecolor{modecol}{RGB}{200,55,55}
\definecolor{meancol}{RGB}{0,90,200}
\definecolor{medcol}{RGB}{160,90,190}
\definecolor{quartcol}{RGB}{0,135,90}
\definecolor{ogivecol}{RGB}{176,74,0}

% ---------------- Boîtes ----------------
\tcbset{boxbase/.style={enhanced,breakable,boxrule=0.9pt,arc=2.5pt,
  left=3mm,right=3mm,top=2.3mm,bottom=2.3mm,
  fonttitle=\bfseries,coltitle=white,toptitle=0.6mm,bottomtitle=0.6mm}}

\newtcolorbox{definition}[1][]{boxbase,colback=defcol!6,colframe=defcol,
  title={\textsc{Définition}\ifx\relax#1\relax\relax\else\textmd{ : \itshape #1}\fi}}
\newtcolorbox{propriete}[1][]{boxbase,colback=thmcol!7,colframe=thmcol,
  title={\textsc{Propriété}\ifx\relax#1\relax\relax\else\textmd{ : \itshape #1}\fi}}
\newtcolorbox{formule}[1][]{boxbase,colback=formulacol!7,colframe=formulacol,
  title={\textsc{Formule}\ifx\relax#1\relax\relax\else\textmd{ : \itshape #1}\fi}}
\newtcolorbox{exemple}[1][]{boxbase,colback=excol!6,colframe=excol,
  title={\textsc{Exemple}\ifx\relax#1\relax\relax\else\textmd{ : \itshape #1}\fi}}
\newtcolorbox{methode}[1][]{boxbase,colback=methcol!8,colframe=methcol,sharp corners,
  title={\textsc{Méthode}\ifx\relax#1\relax\relax\else\textmd{ : \itshape #1}\fi}}
\newtcolorbox{remarque}[1][]{boxbase,colback=remarkcol!7,colframe=remarkcol,
  title={\textsc{Astuce}\ifx\relax#1\relax\relax\else\textmd{ : \itshape #1}\fi}}
\newtcolorbox{attention}[1][]{boxbase,colback=attncol!7,colframe=attncol,
  title={\textsc{Attention}\ifx\relax#1\relax\relax\else\textmd{ : \itshape #1}\fi}}

% Boîte "exercice" et "corrigé" (titre libre passé en argument)
\newtcolorbox{exobox}[1]{boxbase,colback=primary!4,colframe=primary,
  before skip=8pt,after skip=8pt,title={#1}}
\newtcolorbox{corbox}[1]{boxbase,colback=excol!5,colframe=excol!85!black,
  before skip=8pt,after skip=8pt,title={#1}}

% ---------------- En-tête ----------------
\newcommand{\docpart}[1]{\def\thedocpart{#1}}
\docpart{}
\pagestyle{fancy}
\fancyhf{}
\fancyhead[L]{\small\textcolor{primarydark}{\textbf{Fiche 8} — Statistiques descriptives}}
\fancyhead[R]{\small\textcolor{primarydark}{\textsc{\thedocpart}}}
\fancyfoot[C]{\small\thepage}
\renewcommand{\headrulewidth}{0.8pt}
\renewcommand{\headrule}{\hbox to\headwidth{\color{primary}\leaders\hrule height \headrulewidth\hfill}}

\titleformat{\section}{\large\bfseries\color{primarydark}}{\thesection}{0.6em}{}
\titleformat{\subsection}{\normalsize\bfseries\color{primary}}{\thesubsection}{0.6em}{}
\setlist{topsep=2pt,itemsep=1.5pt,parsep=0pt}

% Bandeau de titre du document
\newcommand{\bandeau}[2]{%
\begin{tcolorbox}[boxbase,colback=primary,colframe=primarydark,halign=center,
   before skip=2pt,after skip=12pt]
{\LARGE\bfseries\color{white} #1}\\[3pt]{\normalsize\color{white} #2}
\end{tcolorbox}}

\newcommand{\ecm}{\sigma_m}
\newcommand{\ect}{\sigma}
\newcommand{\med}{M\!e}
\docpart{Difficile — Corrigé — Thème 5}
\begin{document}
\bandeau{Thème 5 — Représentations graphiques}{Niveau Difficile — corrigé}

\begin{corbox}{Corrigé 1 — Du tableau à l'ogive et à la boîte}
\textbf{(a)} La hauteur est une grandeur mesurée regroupée en classes : caractère \textbf{continu},
donc \textbf{histogramme}. Plus grand effectif $16$ : classe modale $\mathbf{[140,160[}$.\\[2pt]
\textbf{(b)} Effectifs cumulés $5,16,32,44,50$, donc $f_c=0{,}10;\,0{,}32;\,0{,}64;\,0{,}88;\,1{,}00$.
Couples pour l'ogive :
\[ (120;0{,}10),\ (140;0{,}32),\ (160;0{,}64),\ (180;0{,}88),\ (200;1{,}00). \]
\textbf{(c)} Rang $\tfrac{N}{2}=25$. Premier cumul $\geqslant25$ : $[140,160[$ ($n_c=32$),
$r=25-16=9$, $n_{\text{classe}}=16$ :
\[ \med=140+\frac{9}{16}\times20=140+11{,}25=\mathbf{151{,}25\ \text{cm}}. \]
Sur l'ogive, on la lirait entre les graduations $140$ et $160$ (là où $f_c$ passe par $0{,}50$).\\[2pt]
\textbf{(d)} Cinq nombres : $x_{\min}=100$ ; $Q_1\approx133{,}6$ ; $\med\approx151{,}3$ ;
$Q_3\approx169{,}2$ ; $x_{\max}=200$.
Écart interquartile : $Q_3-Q_1\approx169{,}2-133{,}6=\mathbf{35{,}6\ \text{cm}}$.
\end{corbox}

\begin{corbox}{Corrigé 2 — Raisonner sur un diagramme circulaire}
\textbf{(a)} Médecine occupe la moitié : $f=0{,}50$, donc $n_{\text{Méd}}=0{,}50\times3000=\mathbf{1500}$.\\[2pt]
\textbf{(b)} Les deux autres sections se partagent le reste : $3000-1500=\mathbf{1500}$ étudiants.\\[2pt]
\textbf{(c)} L'écart cherché est $1500-n_{\text{Sc.bio}}$. Comme Sciences biomédicales $<$ Dentisterie
et que ces deux effectifs somment à $1500$, on a $n_{\text{Sc.bio}}<750$. Donc
\[ 1500-n_{\text{Sc.bio}}\ \in\ ]\,750\,;\,1500\,[. \]
\emph{Élimination :} $600<750$ (exclu) ; $1500$ exigerait $n_{\text{Sc.bio}}=0$ (exclu) ;
$2100>1500$ dépasse même l'effectif de Médecine (exclu). Seule valeur possible : $\mathbf{900}$
(alors Sciences bio $=600$, Dentisterie $=900$, cohérent avec la figure).
\end{corbox}
\end{document}
